% !TEX root = ../../../main.tex
% 所属：第二部分：多变量分析学 / 第2章：多元函数微分学
% 节：2.6 条件极值
% 正文按原文件顺序保存；由 ../chapter2.tex 调用。
\section{条件极值}

目标函数： $f(x_1,\ldots,x_n)$ 与若干约束条件。以三元函数为例，


设目标函数为 $f(x,y,z)$，约束条件为
\[
F(x,y,z)=0,\qquad G(x,y,z)=0.
\]
设 $f,F,G$ 一阶连续可导，且 Jacobi 矩阵
\[
J=\frac{\partial(F,G)}{\partial(x,y,z)}
=\begin{pmatrix}
F_x&F_y&F_z\\
G_x&G_y&G_z
\end{pmatrix}
\]
满秩，即 $\operatorname{rank}J=2$。

我们现在考虑一个必要性的条件.

假设 $(x_0,y_0,z_0)$ 是一个条件极值点。不妨设在此点
\[
\begin{vmatrix}
F_y&F_z\\
G_y&G_z
\end{vmatrix}_{(x_0,y_0,z_0)}\ne0.
\]
由隐函数定理，在 $(x_0,y_0,z_0)$ 的某邻域内，
可由方程组
\[
\begin{cases}
F(x,y,z)=0,\\
G(x,y,z)=0
\end{cases}
\]
解出其中两个未知数，即
\[
y=y(x),\qquad z=z(x).
\]
于是令
\[
\varphi(x)=f(x,y,z)=f(x,y(x),z(x)).
\]
因 $(x_0,y_0,z_0)$ 是条件极值点，$x_0$ 是
$\varphi(x)$ 的一个极值点。由 Fermat 引理，
\[
\begin{aligned}
0=\varphi'(x_0)
&=\left.
\left(
\frac{\partial f}{\partial x}
+\frac{\partial f}{\partial y}\frac{\mathrm dy}{\mathrm dx}
+\frac{\partial f}{\partial z}\frac{\mathrm dz}{\mathrm dx}
\right)\right|_{(x_0,y_0,z_0)}\\
&=\left.
\nabla f(x_0,y_0,z_0)\cdot
\left(
\vect i+\frac{\mathrm dy}{\mathrm dx}\vect j
+\frac{\mathrm dz}{\mathrm dx}\vect k
\right)\right|_{x=x_0}.
\end{aligned}
\]
又由
\[
F(x,y(x),z(x))=0,\qquad G(x,y(x),z(x))=0
\]
可得
\[
\begin{aligned}
\nabla F(x_0,y_0,z_0)\cdot
\left(
\vect i+\frac{\mathrm dy}{\mathrm dx}\vect j
+\frac{\mathrm dz}{\mathrm dx}\vect k
\right)_{x=x_0}&=0,\\
\nabla G(x_0,y_0,z_0)\cdot
\left(
\vect i+\frac{\mathrm dy}{\mathrm dx}\vect j
+\frac{\mathrm dz}{\mathrm dx}\vect k
\right)_{x=x_0}&=0.
\end{aligned}
\]
因此 $\nabla f(x_0,y_0,z_0)$、$\nabla F(x_0,y_0,z_0)$、
$\nabla G(x_0,y_0,z_0)$ 均与同一切向量垂直；
又 $\operatorname{rank}J=2$，
$\nabla F(x_0,y_0,z_0)$ 与 $\nabla G(x_0,y_0,z_0)$ 无关。
从而存在 $\lambda,\mu$，使
\[
\nabla f=\lambda\nabla F+\mu\nabla G,
\qquad
\nabla(f-\lambda F-\mu G)=0.
\]
这就得到了 Lagrange 乘数法。

\begin{proposition}
	构造 Lagrange 函数
\[
L(x,y,z,\lambda,\mu)
=f(x,y,z)-\lambda F(x,y,z)-\mu G(x,y,z).
\]
在 $(x_0,y_0,z_0)$ 处有
\[
\begin{cases}
\displaystyle\frac{\partial L}{\partial x}
=\frac{\partial f}{\partial x}
-\lambda\frac{\partial F}{\partial x}
-\mu\frac{\partial G}{\partial x}=0,\\[6pt]
\displaystyle\frac{\partial L}{\partial y}
=\frac{\partial f}{\partial y}
-\lambda\frac{\partial F}{\partial y}
-\mu\frac{\partial G}{\partial y}=0,\\[6pt]
\displaystyle\frac{\partial L}{\partial z}
=\frac{\partial f}{\partial z}
-\lambda\frac{\partial F}{\partial z}
-\mu\frac{\partial G}{\partial z}=0,\\[6pt]
F(x,y,z)=0,\\
G(x,y,z)=0.
\end{cases}
\]
从而求出 $(x_0,y_0,z_0)$ 以及 $\lambda_0,\mu_0$，
再进一步验证它是否为极值点。
\end{proposition}

条件极值的一般情况：$n$ 个自变量、$m$ 个约束条件，$m<n$。
目标函数为 $f(x_1,\ldots,x_n)$，约束条件为
\[
\begin{cases}
g_1(x_1,\ldots,x_n)=0,\\
\qquad\vdots\\
g_m(x_1,\ldots,x_n)=0.
\end{cases}\tag{*}
\]
所有函数属于 $C^2$，且 Jacobi 矩阵
\[
J=
\begin{pmatrix}
\dfrac{\partial g_1}{\partial x_1}&\cdots&
\dfrac{\partial g_1}{\partial x_n}\\
\vdots&&\vdots\\
\dfrac{\partial g_m}{\partial x_1}&\cdots&
\dfrac{\partial g_m}{\partial x_n}
\end{pmatrix}_{m\times n}
\]
满秩，即 $\operatorname{rank}J=m<n$。

设前 $m$ 个变量由条件 $(*)$ 可以解出唯一解
\[
\begin{cases}
x_1=h_1(x_{m+1},\ldots,x_n),\\
\qquad\vdots\\
x_m=h_m(x_{m+1},\ldots,x_n).
\end{cases}
\]
代入
\[
\varphi(x_{m+1},\ldots,x_n)
=f(h_1,\ldots,h_m,x_{m+1},\ldots,x_n).
\]
将复合函数求导，条件极值点处有
\[
\begin{aligned}
0=\frac{\partial\varphi}{\partial x_{m+1}}
&=\frac{\partial f}{\partial x_1}
\frac{\partial h_1}{\partial x_{m+1}}
+\cdots+
\frac{\partial f}{\partial x_m}
\frac{\partial h_m}{\partial x_{m+1}}
+\frac{\partial f}{\partial x_{m+1}}\\
&=\nabla f\cdot
\underbrace{\left(
\frac{\partial h_1}{\partial x_{m+1}},
\ldots,
\frac{\partial h_m}{\partial x_{m+1}},
1,0,\ldots,0
\right)}_{\vect\eta_1},
\end{aligned}
\]
同理
\[
0=\frac{\partial\varphi}{\partial x_n}
=\nabla f\cdot
\underbrace{\left(
\frac{\partial h_1}{\partial x_n},
\ldots,
\frac{\partial h_m}{\partial x_n},
0,\ldots,1
\right)}_{\vect\eta_{n-m}}.
\]
所以在条件极值点，
$\nabla f$ 与 $\vect\eta_1,\ldots,\vect\eta_{n-m}$ 正交，即
\[
\nabla f\in
\bigl(\operatorname{span}\{\vect\eta_1,\ldots,\vect\eta_{n-m}\}\bigr)^\perp,
\qquad
\dim\operatorname{span}\{\vect\eta_1,\ldots,\vect\eta_{n-m}\}=n-m.
\]
对
\[
g_1(h_1,\ldots,h_m,x_{m+1},\ldots,x_n)=0,\quad
\ldots,\quad
g_m(h_1,\ldots,h_m,x_{m+1},\ldots,x_n)=0
\]
求导，得到
\[
\nabla g_1,\ldots,\nabla g_m
\in
\bigl(\operatorname{span}\{\vect\eta_1,\ldots,\vect\eta_{n-m}\}\bigr)^\perp.
\]
又 $\operatorname{rank}J=m$，故
$\nabla g_1,\ldots,\nabla g_m$ 无关，且
$\nabla f$ 由它们线性表出：
\[
\nabla f=\lambda_1\nabla g_1+\cdots+\lambda_m\nabla g_m,
\qquad
\nabla(f-\lambda_1g_1-\cdots-\lambda_mg_m)=0.
\]
令
\[
L(x_1,\ldots,x_n,\lambda_1,\ldots,\lambda_m)
=f-\lambda_1g_1-\cdots-\lambda_mg_m.
\]
于是必要条件为方程组
\[
\begin{cases}
\displaystyle\frac{\partial L}{\partial x_1}
=\frac{\partial f}{\partial x_1}
-\sum_{k=1}^m\lambda_k\frac{\partial g_k}{\partial x_1}=0,\\
\qquad\vdots\\
\displaystyle\frac{\partial L}{\partial x_n}
=\frac{\partial f}{\partial x_n}
-\sum_{k=1}^m\lambda_k\frac{\partial g_k}{\partial x_n}=0,\\
g_1=\cdots=g_m=0.
\end{cases}\tag{*}
\]

\begin{theorem}
设点 $x_0=(x_1^0,x_2^0,\ldots,x_n^0)$ 及 $m$ 个常数
$\lambda_1,\lambda_2,\ldots,\lambda_m$ 满足方程组 $(*)$。
若方阵
\[
\left(
\frac{\partial^2L}{\partial x_i\partial x_j}
(x_0,\lambda_1,\lambda_2,\ldots,\lambda_m)
\right)_{n\times n}
\]
为正定（负定）矩阵，则 $x_0$ 为满足约束条件的条件极小（大）值点，
因此 $f(x_0)$ 为满足约束条件的条件极小（大）值。
\end{theorem}

\begin{proof}
固定方程组 $(*)$ 给出的乘子
$\lambda_1,\ldots,\lambda_m$，记
\[
A=\left(
\frac{\partial^2L}{\partial x_i\partial x_j}
(x_0,\lambda_1,\ldots,\lambda_m)
\right)_{n\times n}.
\]
由方程组 $(*)$ 可知
$\nabla_xL(x_0,\lambda_1,\ldots,\lambda_m)=0$。
对 $\vect h\to0$，在乘子固定的情况下使用二阶 Taylor 公式，
\[
\begin{aligned}
&L(x_0+\vect h,\lambda_1,\ldots,\lambda_m)
 -L(x_0,\lambda_1,\ldots,\lambda_m)\\
&\qquad=
\frac12\vect h^{\mathsf T}A\vect h
+o\!\left(\lVert\vect h\rVert^2\right).
\end{aligned}
\]

若 $A$ 正定，则存在 $c>0$，使得
$\vect h^{\mathsf T}A\vect h\ge c\lVert\vect h\rVert^2$。
当非零的 $\vect h$ 充分小时，余项的绝对值小于
$\dfrac c4\lVert\vect h\rVert^2$，故
\[
L(x_0+\vect h,\lambda_1,\ldots,\lambda_m)
 -L(x_0,\lambda_1,\ldots,\lambda_m)>0.
\]
对满足约束条件的点 $x_0+\vect h$，有
$g_k(x_0+\vect h)=g_k(x_0)=0$，因而
\[
f(x_0+\vect h)-f(x_0)
=L(x_0+\vect h,\lambda_1,\ldots,\lambda_m)
 -L(x_0,\lambda_1,\ldots,\lambda_m)>0.
\]
所以 $x_0$ 是严格的条件极小值点。
若 $A$ 负定，对 $-L$ 作同样的论证，即得
$f(x_0+\vect h)<f(x_0)$，故 $x_0$ 是严格的条件极大值点。
\end{proof}
